\section{Methods}
\label{sec:nmi-methods}
% Alias label so the reused appendix (SI) cross-references resolve in this variant.
\label{sec:methods}

\paragraph{The scoring system.}
\textsc{AIPR} grades a manuscript with a single frozen pipeline (v6), applying the now-standard
practice of scoring by prompting an LLM against a rubric \citep{liu2023geval,zheng2023judging} to
a new object, the manuscript. It uploads the submitted PDF and runs two passes: a \emph{reviewer}
pass reads the paper and emits five 0--100 dimension scores (novelty, rigor, applicability,
clarity, citation) with justifications and verbatim-anchored highlights, and an \emph{audit} pass
re-reads the paper with a live literature-search tool, checks each highlight, and proposes missing
citations against a real bibliographic index, so a recommended citation is a record the tool
returned rather than a generated string. The overall is a fixed weighted mean,
\begin{equation}
  \text{overall} = \frac{4\,s_{\text{nov}} + 2\,s_{\text{rig}} + 4\,s_{\text{app}}
  + 1\,s_{\text{cla}} + 0.5\,s_{\text{cit}}}{11.5},
  \label{eq:overall}
\end{equation}
with novelty and applicability weighted most and citation least; the model does not emit the
overall directly. The weights were set by internal validation on a separate held-out manuscript
set during product development (no ICLR submissions from any year) and frozen before this study,
not fit to any outcome used here. The system is prompt-only---not fine-tuned on reviews, decisions,
or ratings. Where the baseline prompt is one paragraph of about fifty words, the reviewer pass is
driven by a ${\approx}2{,}200$-word rubric and a $1{,}200$-word output schema and the audit pass
adds a tool loop of up to eight search rounds; the point of the value comparison is that this
instruction footprint, almost two orders of magnitude larger, buys reliability and a grounded
review rather than a higher score.

\paragraph{Configurations and baseline.}
Two configurations run the identical two-pass pipeline and differ only in the model:
\textbf{AIPR (GPT-5.4-mini)} on the cheaper model and \textbf{AIPR (GPT-5.4)} on the frontier
model, the production configuration. The pre-registered baseline \textbf{Direct (GPT-5.4)} runs
the same frontier model on the same PDF with a single one-paragraph prompt and no rubric or audit.
The cheap configuration is effectively unconstrained in volume, the frontier one budgeted, so we
establish the relationship on the large cohort $M$ and confirm on cohort $H$ after verifying the
two agree well enough for the former to proxy the latter (H5).

\paragraph{Cohort and ground truth.}
\label{sec:data}% alias for reused SI cross-references
The primary cohort is ICLR~2026 on OpenReview, chosen because it operates open review---rejected
submissions, reviews, and decisions are public, which the reject-class claim requires---and because
its decisions and reviews were released in January~2026, after the grading model's August~2025
knowledge cutoff, so the outcome could not have been memorized (\S\ref{sec:discussion}). For each
submission we record two ground truths: the decision tier on the ordered ladder reject $<$ poster
$<$ oral (and its binary accept/reject collapse), and the mean reviewer rating, which averages out
individual reviewer noise. We grade the \emph{submitted} PDF a reviewer saw, never the
camera-ready, and record its SHA-256 for provenance. The natural acceptance rate would starve the
accepted tiers, so we sample balanced across tiers and also report metrics re-weighted to the
venue's natural prevalence (\natAcceptRate\% accept). Pre-registered exclusions---desk-rejects,
withdrawals with no reviewer signal, parse failures, and arXiv-twin leakage risks flagged by the
pipeline's self-identity step---never enter a primary metric. Gradings are strictly nested: every
sampled submission is graded on GPT-5.4-mini (cohort $M$, $n=\NminiPrimary$) and a stratified
subset additionally on GPT-5.4 (cohort $H\subseteq M$, $n=\NfullPrimary$), so each frontier
grading has a paired cheap-model score (the bridge, H5). All analysis reads a released two-table
contract (labels per submission, scores per grading run); we redistribute derived metadata and our
own scores under the OpenReview CC~BY~4.0 terms, not manuscript PDFs or full review texts.

\paragraph{Pre-registration and endpoints.}
Hypotheses, the primary metric, and the low-score threshold were fixed before any score was joined
to any outcome (verbatim in the SI). The primary endpoint (H1) is bottom-quintile reject
enrichment, reported as lift over the base rate with Wilson intervals. Confirmatory endpoints are
reject/accept AUROC (H2), a monotone tier gradient (H3), and correlation with the mean reviewer
rating (H4). A gate (H5) requires the cheap and frontier scores to agree at $\rho\ge0.8$ before the
large cohort stands in for the production score (observed \bridgeRhoFull). A pre-registered value
comparison (V1) asks whether the pipeline out-discriminates the baseline; its directional success
criterion was not met and is reported plainly as such.

\paragraph{Statistical analysis.}
Effect sizes lead every claim with 95\% confidence intervals from 4{,}000-resample bias-corrected
and accelerated (BCa) bootstraps; class-conditional statistics such as AUROC use a stratified
resample preserving the observed reject/accept counts. The monotone-trend test is a
10{,}000-permutation Monte Carlo Jonckheere--Terpstra; rates use Wilson intervals; per-dimension
exploratory tests are Benjamini--Hochberg controlled. Secondary covariate, area, weighting, and
disagreement analyses are descriptive and never redefine the primary claim. The entire analysis is
deterministic under a fixed seed and regenerates from the released \texttt{submissions} and
\texttt{gradings} tables. Two estimator choices differ from the registered plan---BCa rather than
percentile intervals, and a whole-cohort bootstrap for the H1 lift CI---both analysis-stage
substitutions of a better-behaved interval that leave every pre-declared threshold unchanged (SI).

\paragraph{Reproducibility and audit posture.}
The de-identified scored dataset, the decision and rating labels, the analysis code, and this
manuscript are openly released at \url{https://github.com/cogeor/aipr-score-validation}
(anchored to the pre-registration tag \texttt{prereg-iclr2026-v2}), so every figure,
statistic, and the entire score-to-outcome analysis reproduces independently from one command. The grading rubric, prompts, and model configuration are proprietary: the
manuscript-to-score function is audited through its outputs---and can be exercised on new
manuscripts via the system---rather than re-implemented, and the prompts and configuration are
available confidentially to editors, reviewers, or auditors. Because the design was pre-registered
before unblinding, this openness lets a single author's validation of a closed instrument be
independently checked rather than taken on trust. The author founded and operates AIPR
(\url{https://aipr.pub}), the system evaluated here, and stands to benefit from a favourable result;
the study was self-funded, and the mitigations are structural---the analysis plan was frozen at a
public tag before any score was joined to an outcome, and the primary comparison includes a baseline
designed to make the pipeline redundant.
